\documentclass[12pt]{article}
\usepackage[utf8]{inputenc}
\usepackage[T1]{fontenc}
\usepackage[margin=1in]{geometry}
\usepackage{setspace}
\usepackage{xcolor}
\usepackage[colorlinks=true, urlcolor=blue!60!black, linkcolor=blue!60!black]{hyperref}
\setlength{\parindent}{0em}
\setlength{\parskip}{0.6em}
\renewcommand{\baselinestretch}{1.0}
\begin{document}
\begin{center}
{\LARGE \textit{Griffin v.\ Breckenridge}, 403 U.S.\ 88 (1971)\par}
\vspace{0.5em}
{\large Case Brief --- Two-Page Edition\par}
\end{center}

\section{Memory Jogger}
Section 1985(3) reaches purely private conspiracies driven by racial or class-based invidious animus, validly grounded in the Thirteenth Amendment and the right of interstate travel (\href{https://supreme.justia.com/cases/federal/us/403/88/}{403 U.S. 88}).

\section{Detailed Case Facts}
In 1966, two white Mississippians, believing that driver R. G. Grady of Memphis was a civil rights worker, blocked his car on a highway near DeKalb, forced Grady and his Black passengers out at gunpoint, and clubbed them about the head. Neither acted under color of law. The victims sued under \S~1985(3), part of the Ku Klux Klan Act of 1871.

\section{Procedural History}
The district court dismissed under \textit{Collins v.\ Hardyman}, which required state action; the Fifth Circuit reluctantly affirmed; the Supreme Court reversed.

\section{Judicial Votes}
9--0 (8--1 on Part V-B). Stewart wrote for the Court; Harlan concurred, declining to join the interstate-travel rationale. No dissents.

\section{Holding}
Section 1985(3) requires no state action but does require racial, or perhaps otherwise class-based, invidiously discriminatory animus. Congress could reach this conspiracy under Section 2 of the Thirteenth Amendment and its power to protect interstate travel. Non-racial animus and Section 5 were reserved.

\section{Analysis of Opinions}
Stewart relied on text (conspiring or going in disguise on the highway, with no color-of-law clause), precedent, the 1871 Act's structure, and legislative history. To avoid creating a general federal tort law, he required class-based animus. Under \textit{Jones v.\ Alfred H.\ Mayer Co.}, Congress may define and remedy the badges and incidents of slavery, and under \textit{Guest} interstate travel is protected against private interference. Harlan found the Thirteenth Amendment sufficient by itself.

\section{Examples: Future Applications}
\textbf{Same-Side:} white supremacists vandalize a Black family's home to drive it out of town. \textbf{Opposite-Side:} anti-union violence (\textit{Carpenters v.\ Scott}) and clinic blockades (\textit{Bray}) lack the required animus. \textbf{Fence-Sitter:} a highway attack on a religious minority, since non-racial animus was reserved.

\section{Critique}
Originalists say the decision stretches the badges and incidents of slavery and grafts an atextual animus requirement onto the statute. Progressives praise the revival of the Klan Act but blame the vague animus test for the later narrowing in \textit{Carpenters} and \textit{Bray}.

\section{Key Quotations}
\begin{quote}
``The language requiring intent to deprive of equal protection, or equal privileges and immunities, means that there must be some racial, or perhaps otherwise class-based, invidiously discriminatory animus behind the conspirators' action. The conspiracy, in other words, must aim at a deprivation of the equal enjoyment of rights secured by the law to all.''
\end{quote}
\noindent --- Justice Stewart, Opinion of the Court, 403 U.S. at 102.

\end{document}
